\begin{notechapter}
\chapter{Surrogate Losses for Classification}\label{ch:surrogate-losses}
\chaptersubtitle{Training on one loss, being judged on another}

\begin{chapterintro}
A classifier is judged by its error rate, the risk under the 0--1 loss, whose
minimizer is the Bayes classifier (\cref{bayes-decision-theory:thm:bayes-classifier}).
Yet nobody trains on the 0--1 loss: it is trained on the cross-entropy, the hinge
loss or another \emph{surrogate}. This chapter asks when that substitution is
safe. The answer has two parts: a surrogate must lead to the Bayes classifier
when minimized over all functions (\emph{calibration}), and its excess risk must
control the excess error rate. You will learn
\begin{itemize}[nosep]
  \item why the 0--1 loss cannot be minimized by gradient methods, and the margin;
  \item the common surrogates (logistic, hinge, exponential, squared) and the
        functions their population risks are minimized by;
  \item classification calibration, and bounds of the excess error rate by the
        excess surrogate risk, proved for the hinge and logistic losses;
  \item proper scoring rules: which losses also recover probabilities;
  \item what label smoothing, class weights and the focal loss change in the
        minimizer.
\end{itemize}
\prerequisites{\cref{ch:bayes-decision-theory,ch:erm}; \cref{ch:cross-entropy,ch:sigmoid-softmax}.}
\end{chapterintro}

\section{The 0--1 loss and the margin}\label{surrogate-losses:sec:margin}

Take two classes, coded $y\in\{-1, +1\}$ in this chapter, and a real-valued score $f(x)$
whose sign is the prediction. The 0--1 loss of the classifier $\sign f$ is
$\bbone[yf(x)\leq0]$ (counting $f(x) = 0$ as an error), so it depends on $f$ and $y$ only through
the \emph{margin}
\[
  m = yf(x),
\]
which is positive for a correct prediction and negative for a wrong one, and whose
size measures confidence. The empirical risk
$\ERisk(f) = \frac1n\sum_i\bbone[y_if(x_i)\leq0]$ of a parametric score $f_\theta$ is piecewise constant in $\theta$:
it changes only when some $f_\theta(x_i)$ crosses zero. Its gradient is zero almost
everywhere and undefined elsewhere, it says nothing about \emph{how} wrong a
prediction is, and minimizing it exactly is NP-hard even for linear classifiers
in general. A \emph{surrogate} replaces the step by a function of the margin that
is convex and decreasing.

\begin{definition}[Margin-based surrogate]
A \emph{margin loss} is a function $\phi:\R\to[0,\infty)$; the loss of score $f$ on $(x, y)$ is $\phi(yf(x))$,
its risk is $\Risk_\phi(f) = \E[\phi(Yf(X))]$, and $\Risk_\phi^* = \inf_f\Risk_\phi(f)$ over all measurable $f$.
\end{definition}

The four classical examples are
\begin{center}
\begin{tabular}{@{}lll@{}}
\toprule
name & $\phi(m)$ & used by \\
\midrule
logistic & $\log_2(1 + e^{-m})$ & logistic regression, neural networks \\
hinge & $\max(0, 1 - m)$ & support vector machines \\
exponential & $e^{-m}$ & AdaBoost \\
squared & $(1 - m)^2$ & least-squares classification \\
\bottomrule
\end{tabular}
\end{center}
The logistic loss with the natural logarithm is the binary cross-entropy of
\cref{sigmoid-softmax:sec:sigmoid}: with $y\in\{-1, 1\}$ and $P(Y = 1\mid x) = \sigma(f(x))$, the probability of
the observed label is $\sigma(yf(x))$ and $-\ln\sigma(m) = \ln(1 + e^{-m})$. Base $2$ only rescales it, and
makes it pass through $(0, 1)$. \Cref{surrogate-losses:fig:losses} draws them.

\begin{figure}[ht]
\centering
\begin{tikzpicture}
\begin{axis}[width=10cm, height=6.5cm, xmin=-2, xmax=2.5, ymin=0, ymax=4.2,
  xlabel={margin $m = yf(x)$}, ylabel={loss}, grid=major, grid style={gray!15},
  legend style={font=\small, draw=none, at={(0.98,0.98)}, anchor=north east}, legend cell align=left]
  \addplot[very thick, black, const plot] coordinates {(-2,1) (0,1) (0,0) (2.5,0)};
  \addlegendentry{0--1}
  \addplot[very thick, bookaccent, domain=-2:2.5, samples=150] {ln(1+exp(-x))/ln(2)};
  \addlegendentry{logistic ($\log_2$)}
  \addplot[very thick, bookexample, dashed, domain=-2:2.5, samples=50] {max(0,1-x)};
  \addlegendentry{hinge}
  \addplot[very thick, bookdef, densely dotted, domain=-2:2.5, samples=150] {exp(-x)};
  \addlegendentry{exponential}
  \addplot[thick, bookpitfall, dashdotted, domain=-2:2.5, samples=150] {(1-x)^2};
  \addlegendentry{squared}
\end{axis}
\end{tikzpicture}
\caption{The 0--1 loss and four convex surrogates, as functions of the margin.
All four pass through $(0, 1)$ and lie above the 0--1 loss. The squared loss also
penalizes margins larger than $1$, i.e.\ predictions that are right with too much
confidence.}
\label{surrogate-losses:fig:losses}
\end{figure}

\paragraph{Upper bounds.} Each of the four is $\geq\bbone[m\leq0]$, so the empirical surrogate
risk bounds the training error. For the cross-entropy written with probabilities,
$-\log p_y$ where $p_y$ is the probability of the true class, this requires the base-2
logarithm: $-\ln p_y<1$ for $\frac1e<p_y<\frac12$, so with natural logarithms a wrong
prediction can cost less than $1$. With $\log_2$ the bound holds for any number of
classes: if the most probable class is wrong, some other class has probability at
least $p_y$, so $p_y\leq\frac12$ and $-\log_2p_y\geq1$. The two versions differ by the factor $\ln2$
and have the same minimizers. Being an upper bound is reassuring but not the
point; what matters is where the minimizer of the surrogate risk lies.

\section{What each surrogate estimates}\label{surrogate-losses:sec:conditional}

By the pointwise principle (\cref{bayes-decision-theory:thm:pointwise}), the surrogate risk is
minimized over all functions by minimizing, at each $x$, the \emph{conditional
surrogate risk}. With $\eta = \eta(x) = P(Y = 1\mid X = x)$ and $a = f(x)$, it is
\[
  C_\eta(a) = \eta\,\phi(a) + (1 - \eta)\,\phi(-a),\qquad
  C^*_\eta = \inf_aC_\eta(a).
\]
Minimizing in $a$ for each loss (\cref{surrogate-losses:exr:1}) gives:
\begin{center}
\begin{tabular}{@{}llll@{}}
\toprule
loss & minimizer $a^*(\eta)$ & minimal risk $C^*_\eta$ & recovers $\eta$? \\
\midrule
logistic ($\ln$) & $\log\frac\eta{1 - \eta}$ & $-\eta\ln\eta - (1 - \eta)\ln(1 - \eta)$ & yes: $\eta = \sigma(a^*)$ \\
exponential & $\frac12\log\frac\eta{1 - \eta}$ & $2\sqrt{\eta(1 - \eta)}$ & yes: $\eta = \sigma(2a^*)$ \\
squared & $2\eta - 1$ & $4\eta(1 - \eta)$ & yes: $\eta = (1 + a^*)/2$ \\
hinge & $\sign(2\eta - 1)$ & $1 - \abs{2\eta - 1}$ & no \\
\bottomrule
\end{tabular}
\end{center}
In all four cases $a^*(\eta)$ has the sign of $2\eta - 1$, so $\sign a^*$ is the Bayes classifier
$\bbone[\eta>\frac12]$ (written in $\pm1$). The first three surrogates estimate a monotone
transformation of the posterior, from which $\eta$ can be read back. The hinge loss
estimates only the Bayes decision: its minimizer jumps from $-1$ to $+1$ at $\eta = \frac12$ and
is the same for $\eta = 0.51$ and $\eta = 0.99$. A support vector machine therefore outputs no
probabilities; the ``probabilities'' some libraries attach to it are fitted
afterwards (Platt scaling). \Cref{surrogate-losses:fig:minimizers} plots the four minimizers.

\begin{figure}[ht]
\centering
\begin{tikzpicture}
\begin{axis}[width=9cm, height=6cm, xmin=0, xmax=1, ymin=-3, ymax=3,
  xlabel={$\eta = P(Y = 1\mid x)$}, ylabel={$a^*(\eta)$}, grid=major, grid style={gray!15},
  legend style={font=\small, draw=none, at={(0.02,0.98)}, anchor=north west}, legend cell align=left]
  \addplot[very thick, bookaccent, domain=0.047:0.953, samples=150] {ln(x/(1-x))};
  \addlegendentry{logistic}
  \addplot[very thick, bookdef, densely dotted, domain=0.0025:0.9975, samples=200] {0.5*ln(x/(1-x))};
  \addlegendentry{exponential}
  \addplot[thick, bookpitfall, dashdotted, domain=0:1, samples=2] {2*x-1};
  \addlegendentry{squared}
  \addplot[very thick, bookexample, dashed, const plot] coordinates {(0,-1) (0.5,-1) (0.5,1) (1,1)};
  \addlegendentry{hinge}
\end{axis}
\end{tikzpicture}
\caption{Minimizers of the conditional surrogate risk as functions of the
posterior. All change sign at $\eta = \frac12$, so all lead to the Bayes classifier; only
the hinge loss forgets the value of $\eta$.}
\label{surrogate-losses:fig:minimizers}
\end{figure}

\section{Calibration and excess-risk bounds}\label{surrogate-losses:sec:calibration}

That $\sign a^*$ is the Bayes classifier is not quite enough. A learning algorithm
never reaches $\Risk^*_\phi$ exactly; we need that \emph{approaching} it forces the error
rate towards $R^*$. This holds if a wrong sign always costs a definite amount of
surrogate risk.

\begin{definition}[Classification calibration]
$\phi$ is \emph{classification-calibrated} if for every $\eta\neq\frac12$,
\[
  \inf\bigl\{C_\eta(a) : a(2\eta - 1)\leq0\bigr\}>C^*_\eta :
\]
scores of the wrong sign are strictly worse than the best score, at every $\eta\neq\frac12$.
\end{definition}

\begin{theorem}[Bartlett, Jordan and McAuliffe, 2006]\label{surrogate-losses:thm:bjm}
A convex margin loss $\phi$ is classification-calibrated if and only if it is
differentiable at $0$ with $\phi'(0)<0$. If $\phi$ is calibrated, there is a nondecreasing
function $\psi$ with $\psi(0) = 0$ and $\psi(\theta)>0$ for $\theta>0$ such that for every measurable $f$,
\[
  \psi\bigl(\Risk(\sign f) - R^*\bigr)\leq\Risk_\phi(f) - \Risk^*_\phi .
\]
\end{theorem}

The four surrogates above are convex with $\phi'(0)<0$, hence calibrated. For two of
them the bound is easy to prove, and the proof shows how the argument works.

\begin{proposition}[Excess-risk bounds]\label{surrogate-losses:prop:bounds}
For every measurable $f$,
\[
  \text{(hinge)}\quad\Risk(\sign f) - R^*\leq\Risk_\phi(f) - \Risk^*_\phi,\qquad
  \text{(logistic, $\ln$)}\quad\Risk(\sign f) - R^*\leq\sqrt{2\bigl(\Risk_\phi(f) - \Risk^*_\phi\bigr)}.
\]
\end{proposition}

\begin{proof}
By \cref{bayes-decision-theory:prop:excess}, the excess error is
$\E\bigl[\abs{2\eta(X) - 1}\,\bbone[W]\bigr]$, where $W$ is the event that $f(X)$ has the wrong sign
($f(X)(2\eta(X) - 1)\leq0$ with $\eta(X)\neq\frac12$). Write $\Delta(x) = C_{\eta(x)}(f(x)) - C^*_{\eta(x)}\geq0$, so that
$\Risk_\phi(f) - \Risk^*_\phi = \E[\Delta(X)]$. It suffices to bound $\abs{2\eta - 1}$ by a function of $\Delta$ on $W$.

On $W$, $a = f(x)$ and the minimizer $a^*$ lie on opposite sides of $0$ (or $a = 0$). Since $C_\eta$ is
convex with minimum at $a^*$, it is monotone between $a^*$ and $a$, so
$C_\eta(a)\geq C_\eta(0) = \phi(0)$, and $\Delta\geq\phi(0) - C^*_\eta$.
\begin{itemize}[nosep]
  \item \emph{Hinge}: $\phi(0) = 1$ and $C^*_\eta = 1 - \abs{2\eta - 1}$, so $\Delta\geq\abs{2\eta - 1}$ on $W$. Hence the
        excess error is at most $\E[\Delta(X)\bbone[W]]\leq\E[\Delta(X)]$.
  \item \emph{Logistic}: $\phi(0) = \ln2$ and $C^*_\eta$ is the binary entropy in nats, so
        $\phi(0) - C^*_\eta = \mathrm{KL}\bigl((\eta, 1 - \eta)\,\Vert\,(\frac12,\frac12)\bigr)$. Pinsker's inequality,
        $\mathrm{KL}(p\Vert q)\geq2\,\mathrm{TV}(p, q)^2$ with $\mathrm{TV} = \abs{\eta - \frac12}$, gives
        $\Delta\geq2(\eta - \frac12)^2 = \frac12(2\eta - 1)^2$, i.e.\ $\abs{2\eta - 1}\leq\sqrt{2\Delta}$ on $W$. Then, by Jensen's
        inequality for the concave square root,
        $\E\bigl[\abs{2\eta - 1}\bbone[W]\bigr]\leq\E\bigl[\sqrt{2\Delta(X)}\bigr]\leq\sqrt{2\E[\Delta(X)]}$.\qedhere
\end{itemize}
\end{proof}

\begin{keyidea}[What the surrogate buys]
Driving the surrogate risk to its minimum over all functions drives the error rate
to the Bayes rate. Combined with the decomposition of \cref{ch:erm}, the excess error
of a trained classifier is controlled by the approximation error of the class
\emph{in surrogate risk} plus the estimation error of the surrogate --- the
quantities that learning theory can bound. The square root for the logistic
loss means the error rate may converge more slowly than the surrogate risk.
\end{keyidea}

\paragraph{Several classes.} With $K$ classes and scores $f_1,\dots,f_K$, the multinomial
logistic loss of \cref{sigmoid-softmax:eq:loss} is calibrated: its population minimizer is
$\softmax(f) = \eta(x)$ (\cref{bayes-decision-theory:prop:logloss}), whose $\argmax$ is the Bayes class. Multiclass
hinge losses are more delicate. The Crammer--Singer loss
$\max\bigl(0, 1 - f_y + \max_{k\neq y}f_k\bigr)$ asks the true score to exceed every other by $1$; the
Weston--Watkins loss $\sum_{k\neq y}\max(0, 1 - f_y + f_k)$ asks it pair by pair. Neither is
calibrated for $K\geq3$: for some conditional distributions --- typically when no class has
conditional probability above $\frac12$ --- their minimizers need not select the most probable
class (T.~Zhang, \emph{Statistical Analysis of Some Multi-Category Large Margin Classification
Methods}, JMLR, 2004; Tewari and Bartlett, \emph{On the Consistency of Multiclass Classification
Methods}, JMLR, 2007). Replacing the hinge by its square restores calibration for the
Weston--Watkins form (Tewari and Bartlett, 2007). Training $K$ binary one-versus-rest classifiers is a third,
different option.

\section{Proper scoring rules}\label{surrogate-losses:sec:proper}

Calibration in the sense above concerns the decision. Often the probability
itself is wanted: to apply a cost matrix (\cref{bayes-decision-theory:thm:costs}), to rank, or to
report risk. A loss $S(q, y)$ of a predicted distribution $q$ is a \emph{proper scoring rule}
if $q\mapsto\E_{Y\sim\eta}[S(q, Y)]$ is minimized at $q = \eta$, and \emph{strictly proper} if $\eta$ is the only
minimizer. Strictly proper losses are the ones whose population minimizer is the
true posterior.
\begin{itemize}[nosep]
  \item The \emph{log score} $S(q, y) = -\log q_y$ is strictly proper: its expectation is
        $H(\eta) + \mathrm{KL}(\eta\Vert q)$ (\cref{cross-entropy:eq:decomposition}).
  \item The \emph{Brier score} $S(q, y) = \sum_k(q_k - \bbone[k = y])^2$ is strictly proper:
        \[
          \E_{Y\sim\eta}\Bigl[\sum_k(q_k - \bbone[k = Y])^2\Bigr] = \sum_k(q_k - \eta_k)^2 + \sum_k\eta_k(1 - \eta_k),
        \]
        the squared distance to $\eta$ plus an irreducible term (\cref{surrogate-losses:exr:2}).
  \item The hinge loss, viewed as a loss on scores, is not proper in any sense: its
        minimizer forgets $\eta$.
\end{itemize}
A model trained with a strictly proper loss estimates probabilities consistently
in the limit, but a large network trained to near-zero training loss is usually
\emph{overconfident} on new data (\cref{sigmoid-softmax:sec:frequencies}). Whether predicted
probabilities match observed frequencies --- among predictions of $0.8$, are $80\%$
correct? --- is measured on held-out data, for instance by the \emph{expected
calibration error} (Naeini et al., 2015), the average gap between confidence and accuracy
over bins of confidence. A simple fix is \emph{temperature scaling} (Guo, Pleiss, Sun and
Weinberger, \emph{On Calibration of Modern Neural Networks}, ICML 2017): divide all
logits by one scalar $T>0$ fitted by minimizing the cross-entropy on a validation set.
It changes the probabilities but not their ranking, hence not the predicted class.

\section{Changing the loss: label smoothing, class weights, focal loss}\label{surrogate-losses:sec:beyond}

Modifications of the cross-entropy are common. The question to ask of each is the
one this chapter has been asking: what does the population risk now estimate?

\paragraph{Label smoothing.} Replace the one-hot target by
$t = (1 - \varepsilon)e_y + \frac\varepsilon K\bbone$ and minimize $-\sum_kt_k\log q_k$. Since the loss is linear in the target,
its expectation over $Y\sim\eta$ is the cross-entropy with target $(1 - \varepsilon)\eta + \frac\varepsilon K\bbone$, so the
minimizer is
\[
  q^* = (1 - \varepsilon)\,\eta + \frac\varepsilon K\bbone .
\]
The posterior is shrunk towards the uniform distribution by a known amount: the
$\argmax$ is unchanged and $\eta$ can be recovered as $(q^* - \frac\varepsilon K)/(1 - \varepsilon)$. No target
probability is $0$ or $1$ any more, so the optimal logits stay finite: the gap between
the largest and smallest is at most $\log\frac{1 - \varepsilon + \varepsilon/K}{\varepsilon/K}$ (\cref{surrogate-losses:exr:3}).

\paragraph{Class weights.} For imbalanced classes the loss of class $k$ is often
multiplied by a weight $w_k$. In the binary case the minimizer becomes
\[
  q^* = \frac{w_1\eta}{w_1\eta + w_0(1 - \eta)},
\]
a distorted posterior whose threshold $q^*>\frac12$ is equivalent to $\eta>\frac{w_0}{w_0 + w_1}$
(\cref{surrogate-losses:exr:4}). Weighting the loss is therefore the same as moving the
decision threshold, which is the cost-sensitive Bayes classifier of
\cref{bayes-decision-theory:thm:costs} with costs proportional to the weights. The output is no
longer an estimate of $\eta$, but it can be corrected by inverting the formula.

\paragraph{Focal loss.} The focal loss (Lin, Goyal, Girshick, He and Doll\'ar, \emph{Focal Loss for Dense Object Detection}, ICCV 2017) multiplies the cross-entropy of
the true class by $(1 - p_y)^\gamma$, $\gamma\geq0$, to reduce the weight of examples that are already
well classified. It is not proper for $\gamma>0$. In the binary case, minimizing
$-\eta(1 - q)^\gamma\log q - (1 - \eta)q^\gamma\log(1 - q)$ numerically for $\eta = 0.9$ gives $q^* = 0.774$ for $\gamma = 1$,
$0.694$ for $\gamma = 2$ and $0.597$ for $\gamma = 5$; for $\eta = 0.6$ and $\gamma = 2$ it gives $0.538$. The minimizer stays
on the same side of $\frac12$ as $\eta$, so the decision is unchanged, but the probabilities are
pulled towards $\frac12$: the model becomes underconfident by design.

\begin{pitfall}[Modified losses estimate modified quantities]
Label smoothing, class weights and the focal loss all change what the trained
model's probabilities mean. That can be the purpose (better decisions on rare
classes, less overconfidence), but the outputs should not then be read as
estimates of $P(Y\mid X)$ without the corresponding correction.
\end{pitfall}

\paragraph{Training with a surrogate.} Whatever the loss, the empirical risk is an
average over examples, and so is its gradient: for the cross-entropy with a
sigmoid output,
\[
  \nabla_\theta\ERisk(f_\theta) = \frac1n\sum_{i=1}^n\bigl(\sigma(f_\theta(x_i)) - y_i\bigr)\nabla_\theta f_\theta(x_i)\qquad(y_i\in\{0,1\}).
\]
A minibatch of $b$ examples drawn uniformly gives an unbiased estimate of this
average, with variance of order $1/b$ (\cref{monte-carlo:sec:plain}); stochastic gradient
descent relies on exactly this. The 0--1 loss gives a gradient of zero, whatever the
batch.

\begin{chaptersummary}
\begin{itemize}[nosep]
  \item The 0--1 loss depends on the margin $m = yf(x)$ through a step; its empirical risk
        is piecewise constant and cannot be minimized by gradients. Convex surrogates
        $\phi(m)$ replace it: logistic, hinge, exponential, squared.
  \item The conditional surrogate risk $\eta\phi(a) + (1 - \eta)\phi(-a)$ is minimized by the logit
        (logistic), half the logit (exponential), $2\eta - 1$ (squared) or $\sign(2\eta - 1)$ (hinge): all
        give the Bayes classifier, only the first three recover $\eta$.
  \item A convex $\phi$ is classification-calibrated iff $\phi'(0)<0$; then the excess error rate is
        controlled by the excess surrogate risk: linearly for the hinge, through a
        square root for the logistic loss.
  \item The multinomial cross-entropy is calibrated for any $K$; the usual multiclass hinge
        losses are not, in general, for $K\geq3$.
  \item The log score and the Brier score are strictly proper; trained networks are
        nevertheless often overconfident, which temperature scaling corrects.
  \item Label smoothing shrinks the estimated posterior towards uniform, class weights
        move the threshold, the focal loss pulls probabilities towards $\frac12$.
\end{itemize}
\end{chaptersummary}

\section*{Exercises}
\addcontentsline{toc}{section}{Exercises}

\begin{exercise}\label{surrogate-losses:exr:1}
Derive the minimizers $a^*(\eta)$ and minimal values $C^*_\eta$ of the conditional risk for the
exponential, squared and hinge losses given in the table of
\cref{surrogate-losses:sec:conditional}.
\solutionlink{surrogate-losses:exr:1}
\end{exercise}

\begin{exercise}\label{surrogate-losses:exr:2}
Prove the decomposition of the expected Brier score in \cref{surrogate-losses:sec:proper}, and
deduce that it is strictly proper.
\solutionlink{surrogate-losses:exr:2}
\end{exercise}

\begin{exercise}\label{surrogate-losses:exr:3}
With label smoothing, show that the minimizer of the expected loss is
$q^* = (1 - \varepsilon)\eta + \frac\varepsilon K\bbone$. If the logits $z$ reach it, show that $\max_kz_k - \min_kz_k\leq\log\frac{1 - \varepsilon + \varepsilon/K}{\varepsilon/K}$, and
evaluate the bound for $K = 1000$ and $\varepsilon = 0.1$.
\solutionlink{surrogate-losses:exr:3}
\end{exercise}

\begin{exercise}\label{surrogate-losses:exr:4}
For the weighted binary cross-entropy $-w_1y\log q - w_0(1 - y)\log(1 - q)$, find the minimizer of the
expected loss for $P(Y = 1) = \eta$, and show that $q^*>\frac12$ iff $\eta>w_0/(w_0 + w_1)$. Which weights
correspond to the cost-sensitive classifier with cost $c_{\mathrm{FN}}$ for a false negative and
$c_{\mathrm{FP}}$ for a false positive?
\solutionlink{surrogate-losses:exr:4}
\end{exercise}

\begin{exercise}\label{surrogate-losses:exr:5}
Show that the squared margin loss $(1 - m)^2$ is calibrated but penalizes some correct
predictions more than some wrong ones, by comparing $m = 3$ with $m = -0.5$. Propose a
modification that avoids this and is still calibrated.
\solutionlink{surrogate-losses:exr:5}
\end{exercise}
\end{notechapter}
